\documentclass[11pt,a4paper]{article}
\usepackage[margin=0.75in]{geometry}
\usepackage{booktabs}
\usepackage{array}
\usepackage{xcolor}
\usepackage{tikz}
\usepackage{enumitem}
\usepackage{fancyhdr}
\usepackage{titlesec}
\usepackage{multicol}
\usepackage{tcolorbox}
\usepackage{fontawesome5}

% Colors
\definecolor{hoopcolor}{RGB}{34, 139, 34}
\definecolor{breakcolor}{RGB}{70, 130, 180}
\definecolor{leavecolor}{RGB}{148, 103, 189}
\definecolor{warmupcolor}{RGB}{255, 140, 0}
\definecolor{aitoncolor}{RGB}{178, 34, 34}

% Section formatting
\titleformat{\section}{\Large\bfseries\color{hoopcolor}}{}{0em}{}[\titlerule]
\titleformat{\subsection}{\large\bfseries}{}{0em}{}
\titlespacing*{\section}{0pt}{12pt}{6pt}
\titlespacing*{\subsection}{0pt}{8pt}{4pt}

% Header/Footer
\pagestyle{fancy}
\fancyhf{}
\lhead{\textbf{Croquet Practice Drills}}
\rhead{Courtside Reference}
\cfoot{\thepage}

% Compact lists
\setlist[itemize]{noitemsep, topsep=2pt, leftmargin=*}
\setlist[enumerate]{noitemsep, topsep=2pt, leftmargin=*}

% Box styles
\tcbset{
    drillbox/.style={colback=gray!5, colframe=gray!50,
        fonttitle=\bfseries, boxrule=0.5pt, arc=2pt,
        left=4pt, right=4pt, top=2pt, bottom=2pt}
}

\begin{document}

\begin{center}
{\Huge\bfseries\color{hoopcolor} Croquet Practice Drills}\\[0.3em]
{\large Courtside Quick Reference}\\[0.5em]
{\small Compiled from Keith Aiton, Wylie, and FVCS Course Materials}
\end{center}

\vspace{0.5em}

%==============================================================================
\section{Warm-Up Routines}
%==============================================================================

\begin{multicols}{2}

\begin{tcolorbox}[drillbox, title={\color{warmupcolor}\faStopwatch~5-Minute Quick}]
\begin{enumerate}
    \item 10 practice swings (no ball)
    \item 5 straight shots at 1m
    \item 5 shots at 2m
\end{enumerate}
\end{tcolorbox}

\begin{tcolorbox}[drillbox, title={\color{warmupcolor}\faStopwatch~10-Minute Full}]
\begin{enumerate}
    \item 10 swings focusing on rhythm
    \item Straight: 5 at 1m, 5 at 2m, 5 at 3m
    \item Clock face: 1 shot each position
    \item 3 roquets at 3m
\end{enumerate}
\end{tcolorbox}

\end{multicols}

%==============================================================================
\section{Fundamental Drills}
%==============================================================================

\begin{tabular}{@{}p{3.5cm}p{7cm}p{4cm}@{}}
\toprule
\textbf{Drill} & \textbf{How To} & \textbf{Goal} \\
\midrule
\textbf{Pendulum Swing} & Swing without ball, smooth backswing \& follow-through & Clean, repeatable motion \\
\textbf{Stop and Stare} & Hit ball, freeze, count to 3 before looking up & Train head stillness \\
\textbf{1m Straight Shot} & Through hoop from 1m, progress when 7/10 & Build confidence \\
\textbf{The Ladder} & Hit to 1m, 2m, 3m, 4m, 5m markers and back & Distance control \\
\bottomrule
\end{tabular}

\vspace{0.5em}

%==============================================================================
\section{Hoop Running}
%==============================================================================

\begin{multicols}{2}

\begin{tcolorbox}[drillbox, title={Clock Face Drill}]
\begin{itemize}
    \item Place markers at 12, 3, 6, 9 o'clock
    \item 2 metres from hoop
    \item Shoot from each position
    \item 3 full rotations
    \item Track weak angles
\end{itemize}
\end{tcolorbox}

\begin{tcolorbox}[drillbox, colframe=aitoncolor!70, title={\color{aitoncolor}Aiton's Approach}]
\textbf{Ideal position:} 1 yard in front, straight on\\[0.3em]
\textit{``Running hoops with control directs your ball to specific positions for advantageous rushes.''}
\end{tcolorbox}

\end{multicols}

%==============================================================================
\section{Roquet \& Rush}
%==============================================================================

\begin{tabular}{@{}p{3.5cm}p{5.5cm}p{5.5cm}@{}}
\toprule
\textbf{Drill} & \textbf{Setup} & \textbf{Progression} \\
\midrule
\textbf{Roquet Practice} & Target ball 2m away, 10 attempts & 2m $\rightarrow$ 3m $\rightarrow$ 4m \\
\textbf{Straight Rush} & Hit directly through centre & Short to long distance \\
\textbf{Cut Rush} & Angle to hit off-centre (like pool) & Gentle cuts $\rightarrow$ sharper \\
\textbf{Roquet Rally} & 2 players, first to 10 roquets wins & Game format practice \\
\bottomrule
\end{tabular}

\vspace{0.5em}

%==============================================================================
\section{Croquet Shot Types}
%==============================================================================

\begin{tabular}{@{}p{2.5cm}cp{9cm}@{}}
\toprule
\textbf{Shot} & \textbf{Ratio} & \textbf{Technique} \\
\midrule
\textbf{Takeoff} & --- & Object stays, striker travels (thin vs thick) \\
\textbf{Stop Shot} & 1:5--1:10 & Strike below centerline, no follow-through \\
\textbf{Drive} & 1:3 & Normal stroke---foundation of break play \\
\textbf{Half Roll} & 1:2 & Downward strike, hands lower, stand forward \\
\textbf{Full Roll} & 1:1 & Both balls equal distance, apply topspin \\
\textbf{Pass Roll} & >1:1 & Steep downward angle---difficult shot \\
\textbf{Split Shot} & varies & Aim halfway between destinations (90\textdegree{} max) \\
\bottomrule
\end{tabular}

\vspace{0.5em}

%==============================================================================
\section{Break Play}
%==============================================================================

\begin{multicols}{2}

\subsection*{Key Positions}
\begin{itemize}
    \item \textbf{Pilot/Reception}: Near current hoop
    \item \textbf{Pioneer}: At your NEXT hoop
    \item \textbf{Pivot}: Mid-court recovery ball
\end{itemize}

\subsection*{Triangle Drill (3-Ball)}
\begin{enumerate}
    \item Set balls in triangle near hoop
    \item Roquet first, take croquet
    \item Roquet second, take croquet
    \item Position and run hoop
\end{enumerate}
\textbf{Goal:} Minimum strokes

\end{multicols}

\begin{tcolorbox}[drillbox, colframe=aitoncolor!70, title={\color{aitoncolor}Aiton on Breaks}]
\begin{itemize}
    \item Rush toward useful balls rather than directly to hoops
    \item Plan pioneer establishment before making first hoop
    \item Identify closest opponent ball to next hoop as intended pilot
    \item ``Shoot at the back ball'' when opponents cluster
\end{itemize}
\end{tcolorbox}

\newpage

%==============================================================================
\section{Leaves}
%==============================================================================

\begin{tabular}{@{}p{2.8cm}p{6.5cm}p{5cm}@{}}
\toprule
\textbf{Leave} & \textbf{Setup} & \textbf{Notes} \\
\midrule
\textbf{NSL} & Opponent NE of H4 (wired), other as H3 pioneer & Bread and butter \\
\textbf{DSL} & Opponents diagonal, partner east as escape & Most forgiving \\
\textbf{MSL} & NSL variant, west ball on H2 wire & Superior pioneer \\
\textbf{Bryant} & One at H4, other 1cm north of H1 & Ultra-aggressive \\
\textbf{Summer} & Ball between H3--4, wired by peg & Most forcing \\
\bottomrule
\end{tabular}

\vspace{0.5em}

\begin{tcolorbox}[drillbox, colframe=leavecolor!70, title={\color{leavecolor}Leave Golden Rules}]
\begin{enumerate}
    \item Final roquet should be your \textbf{partner ball} (organize quality rush)
    \item Complete final hoop off an \textbf{opponent's ball} (partner stays positioned)
    \item Create \textbf{no short shots} for opponent
\end{enumerate}
\end{tcolorbox}

%==============================================================================
\section{Openings (Turns 1--4)}
%==============================================================================

\begin{multicols}{2}

\subsection*{First Ball}
\begin{tabular}{@{}ll@{}}
\textbf{East Boundary} & H4--H5 level (standard) \\
\textbf{Supershot} & Mid-court, aggressive \\
\textbf{Anti-Duffer} & H6 to peg, defensive \\
\end{tabular}

\subsection*{Second Ball (Tices)}
\begin{tabular}{@{}ll@{}}
\textbf{Standard} & West, 8--13yd from C1 \\
\textbf{Duffer} & 2ft N of penultimate \\
\textbf{Micro} & 2--3yd from A baulk \\
\textbf{Super Duper} & Response to wide super \\
\end{tabular}

\end{multicols}

%==============================================================================
\section{Triple Peel Practice}
%==============================================================================

\begin{multicols}{2}

\begin{tcolorbox}[drillbox, colframe=breakcolor!70, title={\color{breakcolor}Standard TP (17 min)}]
Setup from NSL position\\
Full sequence: 4-back, penult, rover peels
\end{tcolorbox}

\begin{tcolorbox}[drillbox, colframe=breakcolor!70, title={\color{breakcolor}Delayed TP (18 min)}]
Start from DSL\\
Opponent in Corner 4
\end{tcolorbox}

\end{multicols}

\begin{tcolorbox}[drillbox, title={Delayed Double Peel---ESSENTIAL}]
\textit{``The Penult peel after 3-back, followed by the Rover peel after Penult is one of the most important things for you to practise. Practise it regularly!''}
\end{tcolorbox}

\subsection*{Recovery Drills (Plan B)}
\begin{tabular}{@{}p{4cm}p{10.5cm}@{}}
\textbf{Jawsed Peel} & Recovery options; intentional jawsing for later completion \\
\textbf{Peel Short} & Options to continue break \\
\textbf{Peel Too Far} & Break adjustment techniques \\
\textbf{Abandon Safely} & Set up good leave when peels fail \\
\end{tabular}

%==============================================================================
\section{Practice Games}
%==============================================================================

\begin{tabular}{@{}p{3cm}ccp{7cm}@{}}
\toprule
\textbf{Game} & \textbf{Players} & \textbf{Time} & \textbf{Focus} \\
\midrule
\textbf{First to Five} & 2--4 & 15--20 min & Same hoop, first to 5 wins---hoop confidence \\
\textbf{Roquet Rally} & 2 & 10 min & Take turns roqueting, first to 10---accuracy \\
\textbf{Golf Croquet} & 2--4 & 20--30 min & First 6 hoops contested---game flow \\
\bottomrule
\end{tabular}

\vspace{1em}

%==============================================================================
\begin{center}
\fbox{\parbox{0.95\textwidth}{
\centering
\textbf{\color{aitoncolor}\large Keith Aiton's Wisdom}\\[0.5em]
\textit{``The gifted amateur practices until he can do it,\\
whereas the professional practices until he cannot NOT do it.''}
}}
\end{center}

\vspace{0.5em}

%==============================================================================
\section*{6-Week Schedule}
%==============================================================================

\begin{tabular}{@{}clp{9cm}@{}}
\toprule
\textbf{Week} & \textbf{Focus} & \textbf{Key Drills} \\
\midrule
1 & Fundamentals & Stance/Grip, 1m Straight Shot, Stop and Stare \\
2 & Accuracy & Increase distance, Clock Face introduction \\
3 & Distance \& Roquet & The Ladder, Roquet Practice \\
4 & Two-Ball Play & Two-Ball Sequence, Roquet Rally \\
5 & Advanced & The Triangle, Croquet Shot Control \\
6 & Strategy & Hoop and Position, Tactical Positioning \\
\bottomrule
\end{tabular}

\vfill

\begin{center}
{\small\textit{Sources: Keith Aiton ``The Basics'', Keith Wylie ``Expert Croquet Tactics'',\\
Fraser Valley Croquet Society, Tame the Triple ePractice Pack}}
\end{center}

\end{document}
